\begin{solapartat}
\punts{0,5} \textbf{a1)} Si el corrent vist des de dalt gira en sentit horari, la cara superior de la bobina serà el pol sud d'un electroimant i la cara inferior de la bobina en serà un pol nord, per tant el camp magnètic generat pel corrent que recorre la bobina estarà dirigit cap avall.

Dos pols nord encarats resultaran en una repulsió entre l'imant i la bobina, per tant la bobina (i la membrana adherida) pujaran.

\punts{0,5} \textbf{a2)} Si el corrent que circula per la bobina és altern, els extrems superior i inferior alternaran la seva polaritat, produint un camp magnètic vertical, que canviarà tota l'estona entre els dos sentits (amunt i avall).

L'alternança entre pol nord i pol sud a la cara inferior de la bobina generarà alternativament repulsió i atracció sobre el conjunt imant+membrana, de manera que el conjunt vibrarà en la direcció vertical amb la mateixa freqüència del corrent altern.
\end{solapartat}

\begin{solapartat}
\punts{0,5} \textbf{b1)} En augmentar el nombre de voltes el camp magnètic augmentarà (de manera proporcional).

\punts{0,5} \textbf{b2)} En augmentar la intensitat del corrent elèctric el camp magnètic augmentarà (de manera proporcional).
\end{solapartat}
